\chapter{Sources and Licensing Notes}

\section*{Primary pedagogical references}

\begin{thebibliography}{9}

\bibitem{activecalculus}
Matthew Boelkins, David Austin, Christina Safranski, and Steven Schlicker.
\emph{Active Calculus: Single Variable}, 2nd ed., 2025.
Creative Commons Attribution-ShareAlike 4.0 International.

\bibitem{thompson}
Silvanus P. Thompson.
\emph{Calculus Made Easy}, 2nd ed., 1914.
Public-domain United States edition distributed by Project Gutenberg.

\bibitem{granville}
William Anthony Granville and Percey F. Smith.
\emph{Elements of the Differential and Integral Calculus}.
Public-domain historical edition.

\bibitem{greenhill}
Alfred George Greenhill.
\emph{Differential and Integral Calculus, with Applications}, 1886.
Public domain.

\end{thebibliography}

\section*{How the sources were used}

This BetterGrades manuscript uses original prose, original organization, original worked examples, and newly composed exercises. The source texts were consulted to verify topic coverage, compare pedagogical sequences, and identify durable classes of calculus problems.

\emph{Active Calculus} was especially useful as a modern reference for moving among numerical, graphical, algebraic, and verbal representations and for emphasizing active student work. Because that book is licensed CC BY-SA 4.0, no exercise or extended passage from it is reproduced verbatim in this manuscript.

\emph{Calculus Made Easy}, Granville and Smith, and Greenhill are public-domain sources in the United States. Historical problem settings drawn from their tradition have been modernized and recomposed rather than mechanically transcribed. Their strongest contribution here is the insistence that calculus answer intelligible questions about motion, rates, approximation, and geometry before becoming a catalogue of symbols.

\section*{Rights-separation policy}

Any future BetterGrades ingestion pipeline should preserve source-level provenance. Public-domain exercises may be stored in a public-domain corpus. Material copied or adapted from CC BY-SA sources must remain separately identified, attributed, and distributed under compatible ShareAlike terms. Original BetterGrades exercises should be marked as original rather than quietly laundering provenance through a database, the sort of administrative miracle that eventually becomes a legal problem.
